% Transcription — fonds Alexandre Grothendieck, shelfmark 15, batch 8 (pages 141-160).
% Established from the facsimile archives/batches/15/batch-08.pdf (a mirror of
% https://grothendieck.umontpellier.fr/15.pdf), per the skill
% transcribe-grothendieck. The batch file carries no cover sheet: PDF page N
% is page 140+N of the folder (the Montpellier footer prints « 141 » ... « 160 »),
% and the archivists' pencil numbers bottom-left agree leaf by leaf wherever
% legible (142, 147, 148, 149, 150, 153, 154, 158, 160 read directly; 145,
% 152, 156, 157, 159 faint but in sequence).
%
% Pass: Opus 5.5 (claude-opus-5-5), 2026-09-23 — first pass, unchecked against the pages by a human.
%
% Thirteen of the twenty pages are transcribed. Absent: pages 145, 152, 155
% and 157, blank (145, 152, 155 show only the ink of the recto through the
% paper); page 159, a leaf of an unrelated typescript (Bourbaki « n° 413 »,
% p. 1-41, on differential operators on a Lie group, scanned upside down,
% nothing in his hand); page 156, a brown cover inscribed in his hand
% « Corps de classes commutatif = Motifs à groupe fondamental motivique
% [connexe] commutatif (= Motifs métabéliens) », whose words become the
% \section heading before page 158, per references/hand.md §5.
%
% What the batch is. Four runs in his hand, none paginated by him.
%   (1) Pages 141-144, the end of §1 begun on page 140 (batch 7): the tori
%   T_K = R_{K/Q} G_m, V_K = Ker N_{K/Q}, their compact quotients. The
%   sentence cut at the foot of page 140 (« E_1 (donc F) est un ») resumes at
%   the head of page 142 (« s-groupe discret de rang r_1+r_2-1 »); page 141,
%   which sits between them, is a separate sheet carrying the explicit form
%   of condition (iii), Corollaire 1 and « (iv) ⟹ (ii) », and is best read
%   after page 142. 142: end of (i) ⟹ (iii), a cancelled (iii) ⟹ (i),
%   (iii) ⟹ (iv) via the character modules Z[I_K] ⊃ N ⊃ Z, with a lattice
%   of fields in the margin (K, K~, K' the fixed field of P, L = K ∩ K').
%   141: Cor. 1 (K « weilien », K_0 = K ∩ R: V_K/V_{K_0} compact, V_{K_0}
%   has no nontrivial compact quotient torus, units of finite index).
%   143: Cor. 2 (a maximal compact quotient V_K^cpt ≃ V_L/V_{L_0}, R =
%   Ker(N_{K/L}) · V_{L_0}); definition S°_K = T_K/R and the diagram (*)
%   1 → V'_K → S°_K → G_m → 1. 144: functoriality (**) under i_{K/K'};
%   Cor. 3 (the projective system of the S°_K is essentially that of the
%   T_K/V_{K_0}, K weilien); Cor. 4 (a unique i : G_m → S°_K with
%   ε i(λ) = λ²). The lower half of 144 is blank.
%   (2) Pages 146-151, headed § 2 (circled, in the margin of 146): CM
%   structures. An ideal 𝔄 of E ⊗_Q K defines a CM structure on E relative
%   to K if the corresponding Σ ⊂ Spec E_{K̄} = Hom(E, K̄) satisfies
%   τ(Σ) = complement for every complex conjugation τ; the field of
%   definition K_0 is weilien, non real. Galois description via I_E × I_K;
%   the condition is symmetric in E, K (147); a long boxed passage struck
%   through. NB (147-148): behaviour under extensions E ⊂ E', K ⊂ K'; a
%   smallest pair (E_0, K_0), both weilien; degree bound binom(2d, d).
%   Bimodules of type CM, the homomorphism Ψ_Σ : T_K → T_E, Ψ_Σ(x) =
%   dét_{t/E}(x_t) = N_{E⊗K/E}(x^{1_Σ}), compatible with N_{K'/K} and
%   i_{E'/E} (diagram, 148). 149: reduction to K = K_0, E = E_0; Ψ_Σ
%   determines Σ; with E fixed, Σ = {σ_1, ..., σ_{d'}}, Ψ_Σ(x) = ∏ σ_i(x),
%   Ψ_Σ = N_{K_r/Q} on K_r, hence Ψ_Σ factors through S°_K. 150: Ψ_Σ i_K =
%   i_{E/Q} (boxed); the image of Ψ_Σ; « Question »: do the Ψ_Σ, Σ running
%   over the CM structures on K, form a faithful family? — answered yes:
%   Ker Ψ_Σ ∩ ⋂ Ker Ψ_{Σ_i} = 1 in S°_K (151). 151: second « Question »,
%   with E quadratic: enough quadratic subfields?, the norm tori
%   N_{K/L} : S°_K → S°_L, Galois group (Z/2Z)^N.
%   (3) Pages 153-154, two separate sheets. 153: K a CM field Galois over
%   Q, group G, ρ complex conjugation; M the G-module of φ : G → Z with
%   φ(ρσ) = -φ(σ), defining a torus T over Q; the subgroup H^1(Q,T)' of
%   classes zero at every finite place: exact sequence with Tate groups
%   Ĥ_0 and Y_∞ (the places at infinity), whence H^1(Q,T)' ≃ combinations
%   of the infinite places with coefficients in F_2 and zero sum. 154, a
%   different subject, cancelled by long pencil diagonals: families of
%   epimorphisms (effective, strict) in a category C, Prop. 0 (stability),
%   Prop. 1 (a set G of objects is generating, resp. strictly generating,
%   iff lim_{G/X} T → X is epimorphic, resp. isomorphic); the proof breaks
%   off at the foot of the page (« tels »), continued presumably elsewhere
%   (page 155 is blank).
%   (4) Pages 158 and 160, after the cover 156, in black ink on older
%   yellowed paper: K a number field, A its ring of integers, the group
%   scheme A^* over Z, units E ⊃ E°; the abelianised Galois group of K as an
%   extension of Pic(A) by A^*(Ẑ)/E°; the connected group scheme G° =
%   A^*/𝔑 attached to K by the classification of abelian motives over K;
%   the extension 0 → G° → G → Γ → 0 (160); fibre functors at a place ξ,
%   M_∘ → Q-spaces with G-action → « réseaux de Hodge », so that a
%   Q-space with G-action acquires, at each place, a Hodge structure on
%   V ⊗ C; a marginal « K~ = corps de classes absolu ».
%
% Notation fixed for the batch, continuing batch 7. T_K, V_K, V'_K, V_{K_0},
% V_L, E_K, E_{1,K} as there; S^{\circ}_{K} for his S with a small circle
% above (« S°_K »); \mathbb{G}_m; \mathbb{U} for the circle group (his
% double-struck U); \mathfrak{A} for the ideal he writes as a script « Vl »
% (§ 2); \Sigma, \Sigma', \Sigma_y (he writes y, γ indiscriminately for the
% base embedding — y here), 1_\Sigma; \Psi_\Sigma; I_E, I_K; K_r for his
% « K_r » (the real subfield); \underline{\underline{A}} for his doubly
% underlined A (a scheme over Z, page 158); \mathfrak{N} for the script
% letter of page 158. The word he writes « weilien » (w-a-i-l-i-e-n, with a
% lead-in stroke) is transcribed weilien throughout; batch 7 prints it
% « Weilien ». His « ss-groupe », « ss-corps », « hom. », « isom. », « tt »,
% « resp. », « cpt », « conj. », « pt » are kept. Plain square brackets in
% the text are his: his own brackets, or words he added in the interline;
% editorial additions are \add{}.
%
% Legibility. Pages 143, 144, 146, 148, 149, 153 read nearly end to end; 141,
% 142, 147, 150, 151 are fast and overwritten, their formulas carry and
% part of the prose does not; the struck boxes of 142 and 147 are
% transcribed only in outline. Page 154 is fast but legible. Nothing on these
% leaves is dated.
%
% The dating is the inventory's own for the folder, copied verbatim. Its
% group, [10-18] « Théorie des motifs », is dated [à partir de 1958-à partir
% de 1982].
\documentclass[11pt,a4paper]{article}
% Le préambule commun à toutes les transcriptions du fonds Grothendieck.
%
% Il tient en une idée : **l'incertitude se compose, elle ne se résout pas.**
% Une transcription qui lisse un mot douteux a détruit la seule chose pour
% laquelle on la faisait — un lecteur doit pouvoir savoir, à chaque endroit,
% ce qui était sur le feuillet et ce qui a été ajouté. Les six macros qui
% suivent sont donc l'appareil critique, et non de la décoration.
%
% Les mêmes commandes sont reconnues par `scripts/render.mjs`, qui produit la
% vue de lecture du site. Une macro ajoutée ici doit l'être aussi là-bas,
% faute de quoi le rendu échouera bruyamment — ce qui est voulu : un rendu
% silencieusement faux est pire qu'un rendu absent.

\ProvidesPackage{grothendieck}[2026/08/08 Transcriptions du fonds Grothendieck]

\usepackage{amsmath,amssymb,amsthm}
\usepackage{mathrsfs}
\usepackage{stmaryrd}
% Les diagrammes commutatifs sont partout dans ce fonds : c'est la langue
% même de la géométrie algébrique de Grothendieck, et une transcription qui
% les rendrait en prose ne transcrirait rien.
\usepackage{tikz-cd}
% Français par défaut — c'est la langue des feuillets — mais l'anglais est
% chargé aussi : la traduction et le résumé sont des documents anglais, et la
% typographie française (espaces avant les deux-points) y serait une faute.
% Ces documents basculent par \selectlanguage{english} en tête de corps.
\usepackage[english,french]{babel}
\usepackage{fontspec}
\usepackage{geometry}
\geometry{a4paper,margin=25mm,marginparwidth=28mm}
\usepackage{marginnote}
\usepackage{xcolor}
% ulem, not soul. `soul`'s \st and \ul are built on a character-by-character
% scan that XeLaTeX's font machinery breaks: the document fails with an
% undefined control sequence rather than degrading. ulem does the same two jobs
% and is Unicode-engine safe. `normalem` stops it hijacking \emph, which the
% transcriptions use for Grothendieck's own emphasis.
\usepackage[normalem]{ulem}
\usepackage{hyperref}
\hypersetup{colorlinks=true,linkcolor=black,urlcolor=black,pdfborder={0 0 0}}

% --- Métadonnées du lot -----------------------------------------------------
% Reprises telles quelles par le script de rendu, qui les affiche en tête de
% la vue de lecture. Elles fixent ce que le fichier prétend couvrir : sans
% elles, une transcription flotte sans point d'attache dans un fonds de seize
% mille feuillets.

\newcommand{\folder}[1]{\gdef\grothFolder{#1}}
\newcommand{\batch}[1]{\gdef\grothBatch{#1}}
\newcommand{\leaves}[2]{\gdef\grothFirst{#1}\gdef\grothLast{#2}}
\newcommand{\foldertitle}[1]{\gdef\grothFoldertitle{#1}}
\gdef\grothFolder{?}\gdef\grothBatch{?}\gdef\grothFirst{?}\gdef\grothLast{?}\gdef\grothFoldertitle{}

% --- Le résumé --------------------------------------------------------------
%
% Il ouvre la lecture modernisée, sous le titre « Résumé », et il est composé
% en petit. Ce n'est pas un ornement : c'est le seul passage du document qui
% ne soit pas des mathématiques, et un lecteur doit pouvoir le reconnaître
% avant de l'avoir lu — puis le sauter s'il connaît déjà le sujet, ou s'y
% arrêter s'il ne le connaît pas. Le filet qui le suit marque la reprise du
% corps.

\newenvironment{resume}
  {\small\setlength{\parskip}{0.45em}}
  {\par\medskip\hrule\medskip}

% --- Mots-clés ---------------------------------------------------------------
%
% En anglais, en clôture du résumé de la lecture modernisée : le vocabulaire
% moderne sous lequel chercher ce que les feuillets construisent. Le manifeste
% du site (`npm run manifest`) les extrait pour étiqueter la cote entière —
% cette ligne est la source unique des tags, il n'y a pas de fichier à côté.
\newcommand{\keywords}[1]{\par\smallskip\noindent
  {\footnotesize\textsc{Keywords} --- \textit{#1}}\par}

% --- L'appareil critique ----------------------------------------------------

% Le numéro de feuillet, en marge. C'est l'ancre qui permet de régler un
% désaccord sur une formule en regardant *un* feuillet plutôt qu'une chemise
% de six cents. Le site s'en sert aussi pour tourner les pages du fac-similé
% quand on fait défiler la transcription.
\newcommand{\leaf}[1]{%
  \marginnote{\footnotesize\color{gray!70}\textsf{#1}}%
  \label{leaf:#1}\ignorespaces}

% Illisible : on ne devine pas. Un mot inventé qui se lit comme les autres est
% le pire résultat possible.
\newcommand{\ill}{\textcolor{red!60!black}{[\,\ldots\,]}}

% Lecture douteuse : le mot est donné, et signalé comme incertain.
\newcommand{\uncertain}[1]{\uline{#1}}

% Ajout de l'éditeur — une lettre manquante, un mot sous-entendu.
\newcommand{\add}[1]{\textcolor{blue!50!black}{[#1]}}

% Biffé par Grothendieck lui-même. Ce qu'il a rayé fait partie du document :
% une rature dit souvent où la pensée a changé de direction.
\newcommand{\struck}[1]{\sout{#1}}

% Note du transcripteur, sur l'état matériel du feuillet.
\newcommand{\note}[1]{\par\smallskip{\small\color{gray!60!black}\textit{[#1]}}\par\smallskip}

% Note marginale de Grothendieck, distincte du corps du texte.
\newcommand{\marginal}[1]{\par\smallskip\noindent\hspace{1em}%
  {\small\color{gray!60!black}#1}\par\smallskip}

% --- Titre ------------------------------------------------------------------

\AtBeginDocument{%
  \begin{center}
    {\large\bfseries Fonds Alexandre Grothendieck --- cote n\textsuperscript{o}~\grothFolder}\\[2pt]
    {\itshape\grothFoldertitle}\\[4pt]
    {\small feuillets \grothFirst--\grothLast\ (lot \grothBatch)}\\[2pt]
    {\footnotesize\color{gray!60!black}Transcription établie d'après le fac-similé numérisé
     par l'Université de Montpellier.\\
     Les lectures incertaines sont soulignées, les illisibles notées [\,\ldots\,],
     les ajouts entre crochets.}
  \end{center}
  \vspace{1em}}

\endinput


\folder{15}
\batch{8}
\pages{141}{160}
\dating{1965-[vers 1977]}
\watermark{Édition de démonstration}
\foldertitle{Théorie arithmétique. Théorie de Galois des motifs (1965) : notes manuscites (s.d.), tapuscrit (s.d.).}

\begin{document}

\page{141}
\note{la phrase interrompue au bas du feuillet 140 (lot 7) ne reprend pas
ici mais en tête du feuillet 142 ; ce feuillet-ci, qui explicite la
condition (iii) et donne le Corollaire 1, se lit à la suite du feuillet
142}

Pour expliciter la condition (iii), il suffit donc : expliciter les tores
\uncertain{compacts} quotients du \struck{$T$} $V_{L}$, et prouver que ce sont
les quotients de $V_{L}/V_{L_{0}}$. \struck{i.e.\ \ill{}}
\struck{$V_{L_{0}}$ n'est \uncertain{majoré} par des [tores] quotients
compacts}, donc \note{en interligne, au-dessus de la fin de la ligne biffée :
« sauf le premier membre » (lecture douteuse)}

\struck{Le plus grand quotient compact de $V_{L}$ est majoré par
$V_{L}/V_{L_{0}}$. Il reste : prouver que ce dernier est compact.
\struck{Posons} C'est trivial si $L = L_{0}$. Sinon, (faisant $L = K$ pour
simplifier les notations), considérons \struck{i.e.\ que \ill{} $V_{L}/V_{L}$}
i.e.\ \ill{} [la première partie] prouver des \ill{}}
\[
  \struck{L_{\mathbb{R}} \simeq \mathbb{R}^{r}} \qquad
  L_{0\mathbb{R}} = \ill{}
\]
\note{tout ce paragraphe est barré de quatre longs traits obliques et, en
partie, encadré ; la formule finale est laissée inachevée}

\emph{Corollaire 1}. Supposons $K$ weilien, soit $K_{0} = K \cap \mathbb{R}$.
Alors $V_{K}/V_{K_{0}}$ est compact, et $V_{K_{0}}$ n'a d'autre [tore]
quotient compact que le tore unité. \struck{L'image de $E_{K}$ dans} [De plus
$E_{K} \cap V_{K_{0}}(\mathbb{Q})$ est d'indice fini dans $E_{K}$, i.e.\ (pour
la \uncertain{première} assertion)] l'image de $E_{K}$ dans
$V_{K}/V_{K_{0}}$ est finie.
\note{« weilien » : le mot s'écrit w-a-i-l-i-e-n, précédé d'un trait
d'attaque ; c'est le mot que le lot 7 (feuillet 140) transcrit « Weilien »,
pour le plus grand sous-corps $L$ de $K$ tel que $L_{0} = L \cap \mathbb{R}$
soit d'indice $1$ ou $2$ dans $L$. Il est transcrit ainsi dans tout le lot}

Il suffit de prouver l'\ill{} analogue \ill{} puis \struck{\ill{}} [dans
$V_{L_{0}}$ et \ill{}] de la base $\mathbb{Q} \to \mathbb{R}$. Alors
$K_{0\mathbb{R}} \simeq \mathbb{R}^{r}$, \struck{\ill{}} $V_{L_{0}}$ est un
[tore] déployé, donc tout quotient est déployé, donc compact seulement s'il
est nul. D'autre part $V_{K}/V_{K_{0}}$ est compact : c'est trivial si $K =
K_{0}$, sinon on a $K_{\mathbb{R}} \simeq \mathbb{C}^{r}$, et
\[
  V_{K}/V_{K_{0}} \simeq T_{K}/T_{K_{0}} \simeq
  \Bigl[\Bigl(\prod_{\mathbb{C}/\mathbb{R}} \mathbb{G}_{m}\Bigr) \big/
  \mathbb{G}_{m}\Bigr]^{r} \simeq \mathbb{U}^{r}
\]
est compact. \note{devant le crochet, un premier signe $\prod$ est biffé}

La dernière \uncertain{inclusion} \struck{\ill{}} signifie que $E_{K_{0}}$ est
d'indice fini dans $E_{K}$. Or les rangs sont respectivement, si
$[K_{0} : \uncertain{K}] = r$, $r-1$ et $r-1$, OK.
\note{on attendrait $[K_{0} : \mathbb{Q}] = r$ ; la lettre lue est un $K$}

(iv) $\Longrightarrow$ (ii) Résulte immédiatement du fait que $N_{L/K}$
\uncertain{applique} $E_{1K}$ dans $E_{1L}$, et de la dernière assertion du
corollaire.

\page{142}
\note{la page reprend la phrase interrompue au bas du feuillet 140 (lot 7),
« et $E_{1}$ (donc $F$) est un »}
s-groupe \emph{discret} de rang $r_{1} + r_{2} - 1$, de sorte que le quotient
$V(\mathbb{R})/F$ est compact. Il en est \uncertain{a fortiori} de même de
$V(\mathbb{R})/\overline{F}$, donc de $V'_{\mathbb{R}}$ qu'il majore.

\struck{(iii) $\Longrightarrow$ (i), car si $V'_{\mathbb{R}}$ est compact
i.e.\ $V'(\mathbb{R})$ compact, considérant $V'(\mathbb{R})$ comme quotient de
$\mathbb{R}^{*+(r_{1}+r_{2}-1)} \times (\mathbb{Z}/2\mathbb{Z})^{r_{1}}
\times \mathbb{U}^{r_{2}}$, on voit que ce sera un quotient de
$(\mathbb{Z}/2\mathbb{Z})^{r_{1}} \times \mathbb{U}^{r_{2}}$, or l'image de
$E_{1}$ dans ce dernier est déjà finie [l'image \ill{}]}
\note{paragraphe encadré, barré de traits obliques et d'un trait
horizontal}

(iii) $\Longrightarrow$ (iv) Si $V'_{K}$ est un tore compact, les groupes des
car.\ à \struck{\ill{}} valeurs dans \struck{\ill{}} $\overline{\mathbb{Q}}$
[$\subset \mathbb{C}$], regardés comme modules sous
\[
  G = \mathrm{Gal}(\widetilde{K}/K)
\]
($\widetilde{K}$ \struck{$\subset \mathbb{Q}$} une ext.\ de $K$ finie
galoisienne sur $\mathbb{Q}$) et tels que $\tau \in G$ [la conjugaison
complexe] y opère par $-\mathrm{id}_{M}$, \ill{} \ill{}, \ill{} $G'$ \ill{}
[donc est central, donc $M$ est un $G' = G/P$-module].
\struck{Donc si $G' = G/H$ est le quotient de $G$ qui \ill{} \ill{}
[quotient de Galois de la plus grande \ill{} ext.\ weilienne] \ill{} une
extension $K' \subset \widetilde{K}$ \ill{} $R \supset
\mathrm{Ker}(N_{K/L} | V_{K})$ i.e.\ que $V_{L}$ est un quotient de $V_{K}$.
Prouvons d'abord que \ill{} considérons \ill{} $\tau' = \mathrm{Im}(\tau
\mapsto G')$ est dans le centre de $G'$, donc $K'$ est weilien.}
\note{passage très surchargé, barré de plusieurs traits et en partie encadré ;
seul le squelette en est rendu}

$I_{K}$ [$= G/Q$] espace homogène sous $G$ des hom.\ de $K$ dans
$\overline{\mathbb{Q}}$ ; $\mathbb{Z}[I_{K}]$ [groupe libre engendré par
$I$] groupe des caractères de $T_{K}$ à valeurs dans $\overline{\mathbb{Q}}$,
comme $G$-module. \struck{$N \subset \mathbb{Z}[I]$ tel que}
$\mathbb{Z} \subset \mathbb{Z}[I_{K}]$ fonctions constantes sur $I$.

\begin{tikzcd}
\mathbb{Z} \arrow[r, hook] & N \arrow[r, hook] & \mathbb{Z}[I_{K}] \\
\mathbb{Z} \arrow[u, "\wr" description, no head] \arrow[rr, hook] & & \mathbb{Z}[I_{L}] \arrow[u, hook]
\end{tikzcd}

$N$ sous $G$-module, tel que $\tau$ opère sur $N/\mathbb{Z}$ par
$-\mathrm{id}$, donc $P$ y opère trivialement ; donc $H$ opère trivialement
sur $N$, donc $N \subset \mathbb{Z}[I_{K}]^{P} = \mathbb{Z}[P \backslash
I_{K}]$. \struck{\ill{}} Or $P \backslash I_{K} = I_{L}$, donc on trouve $N
\subset \mathbb{Z}[I_{L}]$, ce qui signifie que le quotient $V'$ de $V_{K}$
est majoré par $V_{L}$.

\marginal{treillis de corps dans la marge gauche :}

\begin{tikzcd}[column sep=small, row sep=small, nodes={font=\scriptsize}]
K \arrow[r, no head, "Q"] \arrow[d, no head, "H"'] & \widetilde{K} \arrow[d, no head, "P"] \\
L = K \cap K' \arrow[r, no head] \arrow[d, no head] & K' \\
\mathbb{Q} \arrow[ur, no head, "G' = G/P"'] \arrow[uur, no head, bend right=40, "G"'] &
\end{tikzcd}

\marginal{$I_{K} = G/Q$, $I_{L} = G/PQ = P \backslash I_{K}$ ; $\tau \in G$
d'ordre $2$}
\note{dans la figure, les deux arcs partent de $\mathbb{Q}$ et aboutissent à
$K'$ et à $\widetilde{K}$ ; au-dessus de $K$, un $\widetilde{K}$ biffé}

\page{143}
\emph{Corollaire 2}. Parmi les $V'$ satisfaisant les conditions de la
proposition, il y en a un maximal [$V_{K}^{\mathrm{cpt}}$
\struck{\ill{} $V^{\mathrm{cpt}}$}], correspondant à un $R$ minimal. Le $R$
correspondant est connu. Avec les notations de (iv), [$N_{L/K}$ induit un
isom.\ $V_{K}^{\mathrm{cpt}} \simeq V_{L}/V_{L_{0}} = V_{L}^{\mathrm{cpt}}$,
le $R$ correspondant] c'est aussi le produit
$[\mathrm{Ker}(N_{K/L})] \cdot V_{L_{0}}$. \struck{et l'hom\ill{}}

Seule la dernière assertion demande encore une vérification. Comme
$N_{K/L}(\lambda) = \lambda^{r}$ si $\lambda \in L$, on voit que le $R$
cherché contient $V_{L_{0}}$, et aussi bien sûr $\mathrm{Ker}(N_{K/L})$.
\struck{Ou bien} donc $R$ contient le produit $R'$. \struck{Alors} Comme
\uncertain{bien} $\lambda \mapsto \lambda^{r}$ de $V_{L_{0}}$ dans lui-même
est un épimorphisme, on trouve que $N_{L/K}$ induit un isom.
\[
  V_{K} / (\mathrm{Ker}(N_{K/L}) \cdot V_{L_{0}}) \simeq
  V_{L} / N_{L/K}(V_{L_{0}}) \simeq V_{L}/V_{L_{0}} ,
\]
ce qui établit le corollaire 2.

On désigne par $S^{\circ}_{K}$ le quotient $T_{K}/R$, donc on a une suite
exacte
\[
  1 \longrightarrow V'_{K} \longrightarrow S^{\circ}_{K}
  \xrightarrow{\;\varepsilon\;} \mathbb{G}_{m} \longrightarrow 1
\]
déduite d'une suite exacte

\begin{tikzcd}
1 \arrow[r] & V_{K} \arrow[r] \arrow[d, "\mathrm{can}"] & T_{K} \arrow[r, "N_{K}"] \arrow[d, "\mathrm{can}"] & \mathbb{G}_{m} \arrow[r] \arrow[d, no head] & 1 \\
1 \arrow[r] & V'_{K} \arrow[r] & S^{\circ}_{K} \arrow[r] & \mathbb{G}_{m} \arrow[r] & 1
\end{tikzcd}

\marginal{$(*)$}
\note{le trait vertical de droite est un double trait d'égalité}

Pour $K$ variable, $S^{\circ}_{K}$ est contravariant en $K$ pour
$N_{K'/K}$, ainsi que le diagramme $(*)$ [\struck{\ill{}} la fonction du
plus gd quotient cpct est fonctorielle en le tore $T$ sur $\mathbb{Q}$]. Il
est aussi covariant \struck{\ill{}} par $i_{K,K'}$ [\uncertain{même}
raison], mais bien sûr cette injection n'est plus compatible avec
$\varepsilon$, on a de façon précise, pour $K \subset K'$

\page{144}

\begin{tikzcd}
1 \arrow[r] & V'_{K} \arrow[r] \arrow[d] & S^{\circ}_{K} \arrow[r] \arrow[d, "i_{K/K'}"] & \mathbb{G}_{m} \arrow[r] \arrow[d, "\lambda \mapsto \lambda^{r}"] & 1 \\
1 \arrow[r] & V'_{K'} \arrow[r] & S^{\circ}_{K'} \arrow[r] & \mathbb{G}_{m} \arrow[r] & 1
\end{tikzcd}

\marginal{$(**)$}
$r = [K' : K]$

On trouve :

\emph{Corollaire 3}. Le système projectif des $S^{\circ}_{K}$, pour $K$
sous-extension finie de $\overline{\mathbb{Q}}$, est \uncertain{essentiellement}
\ill{} (\ill{} \ill{}), \uncertain{équivalent} au syst.\ projectif des
$S^{\circ}_{K}$ pour $K \subset \overline{\mathbb{Q}}$, $K$ weilien (et
galoisien \ill{}), on \uncertain{trouve} un syst.\ projectif des
\struck{$T_{K}$} $T_{K}/V_{K_{0}}$, où $K \subset \overline{\mathbb{Q}}$ est
weilien et $K_{0} = K \cap \mathbb{R}$.

\emph{Corollaire 4}. Soit $K$ ext.\ finie de $\mathbb{Q}$. Il existe un et un
seul hom
\[
  \mathbb{G}_{m} \xrightarrow{\;i\;} S^{\circ}_{K}
\]
tel que $\varepsilon i(\lambda) = \lambda^{2}$.

Comme $\mathrm{Ker}\,\varepsilon$ est compact, et les hom $\mathbb{G}_{m} \to
\mathrm{Ker}\,\varepsilon$ sont nuls, l'unicité est claire. Pour
l'existence, on peut donc supposer que $K$ est weilien. Si \ill{} [i.e.\
$K_{0} = K$] $K \subset \mathbb{R}$, on a $S^{\circ}_{K}
\xrightarrow[\sim]{\;\varepsilon\;} \mathbb{G}_{m}$, et il suffit de prendre
$i(\lambda) = \lambda^{2}$. Si $K_{0} \neq K$, alors on prend pour $i =
i_{K_{0},K} : S^{\circ}_{K_{0}} = \mathbb{G}_{m} \to S^{\circ}_{K}$, en
utilisant la commutativité de $(**)$ (où $r = 2$).

\subsection*{§ 2}

\note{« § 2 », entouré, dans la marge gauche du feuillet 146}

\page{146}
\marginal{\ill{} \ill{} $\Sigma$ de $E \otimes K$}
Soient $E, K$ deux extensions de $\mathbb{Q}$, \struck{$\Sigma$}
$\mathfrak{A}$ un idéal de $E \otimes_{\mathbb{Q}} K$ [$= E_{K}$],
$\overline{K}$ une clôture algébrique de $K$. On dit que $\mathfrak{A}$
définit une CM-structure sur $E$, relativement à $K$, si l'idéal
$\mathfrak{A}'$ de $E \otimes_{\mathbb{Q}} \overline{K} = E_{K}
\otimes_{K} \overline{K}$ engendré par $\mathfrak{A}$ correspond à un
sous-ensemble $\Sigma$ de l'ensemble \struck{\ill{}}
\[
  \mathrm{Spec}\, E_{\overline{K}} \simeq \mathrm{Hom}_{\mathrm{alg}}(E,
  \overline{K})
\]
qui est tel que pour tout élément $\tau$ de $G = \mathrm{Gal}(\overline{K}/K)$
qui est une conjugaison, on ait $\tau(\Sigma) = (\mathrm{Spec}\,
E_{\overline{K}}) - \Sigma$. Notez que ceci implique que pour $g \in G$, on a
aussi $g \tau g^{-1}(\Sigma) = (\mathrm{Spec}\, E_{\overline{K}}) - \Sigma =
\tau \Sigma$, donc $\tau^{-1} g \tau g^{-1}(\Sigma) = \Sigma$, donc $\Sigma$
est invariant \struck{\ill{}} par le s-groupe \struck{\ill{}} fermé invariant
$H$ engendré par les $\tau^{-1} g \tau g^{-1}$. Cela signifie que, si $K_{0}
\subset K$ est le corps de définition de l'idéal $\mathfrak{A}$ de $E_{K}$,
le corps $K_{0}$ [$= K_{\Sigma}$] est \emph{weilien} ; \struck{et} il est
d'ailleurs weilien \emph{non réel}, car $\tau(\Sigma) \neq \Sigma$ [puisque
sinon on aurait $\mathrm{Spec}\, E_{\overline{K}} = \emptyset$, ce qui n'est
pas]. \struck{Alors le groupe de \ill{} $\overline{\mathbb{Q}}$.}

Soit $\overline{\mathbb{Q}}$ une clôture \emph{algébrique} de $\mathbb{Q}$
\struck{Supposons maintenant également $K$ fini sur $\mathbb{Q}$}, [de]
groupe de Galois $G$. Alors $E, K$ correspondent à des espaces homogènes
compacts sous $G$, soient $I_{E}, I_{K}$. La donnée de $\mathfrak{A}$
correspond à celle d'une partie fermée de $I_{E} \times I_{K}$, invariante
par $G$. Choisissons un plongement $K \subset \overline{\mathbb{Q}}$, donc
$I_{K} \simeq G/Q$, $Q$ \struck{le} sous-groupe de $G$ qui laisse fixe les
éléments de $K$. \struck{L'hypothèse Remplaçant} Prenant $K$ pour
$\overline{K}$ [$\overline{\mathbb{Q}} =$] (pour \ill{} des \ill{}
$\overline{\mathbb{Q}}$), et $\Sigma$ \ill{} $I_{K}$ pour son pt origine $y$
\struck{\ill{}} [le plongement $y$ de $K$ dans $\overline{\mathbb{Q}}$], soit
$\Sigma' = (I_{E} \times I_{K}) - \Sigma$, d'où $\Sigma'_{y}$. [Écrire]
$\Sigma_{y} = \Sigma \cap (I_{E} \times y)$. L'hypothèse est que pour tte
conjugaison $\tau \in G$,
\struck{$(\tau \times 1)(\Sigma \cap (I_{E} \times y)) = (I_{E} - \Sigma)$}
\[
  (\tau \times 1)\Sigma_{y} = \Sigma'_{y}
\]
\struck{$(\tau \times 1)\Sigma_{y} = I_{E} - \Sigma_{y}$}
i.e.\ \struck{$(\tau \times 1)\Sigma_{y} \subset I_{E} - \Sigma$,
$(\tau \times 1)\Sigma'_{y} \subset I_{E} - \Sigma_{y}$}
\struck{pour \ill{} écrire aussi}
\note{le bas de la page est une suite d'essais biffés de la même
condition}

\page{147}
Il est évident que la condition envisagée ne dépend pas du choix de
l'immersion $K \hookrightarrow \overline{\mathbb{Q}}$, \struck{\ill{}} i.e.\
du choix de $y$, et peut donc s'écrire
\[
  (\tau \times 1)\Sigma = \Sigma' \qquad \text{pour tte conjugaison } \tau
\]
\struck{$(\ill{})$}
Je dis qu'elle équivaut aussi à
\[
  (1 \times \tau)\Sigma \uncertain{\subset} \Sigma' \qquad
  \text{pour tte conjugaison } \tau
\]
\note{le signe entre $\Sigma$ et $\Sigma'$ dans la seconde formule ressemble
à $\subset$ ; dans la première, un $=$ appuyé recouvre peut-être un autre
signe}
i.e.,

En effet, comme $(\tau \times \tau)\Sigma = \Sigma$, la relation
\uncertain{envisagée} équivaut à
\[
  (\tau \times 1)(\tau \times \tau)\Sigma = \Sigma' \quad \text{i.e.} \quad
  (1 \times \tau)\Sigma = \Sigma' .
\]
\emph{Donc la condition envisagée est symétrique en $E, K$.}

\struck{Considérons maintenant le sous-corps $K_{0}$ de $K$, corps des
invariants du sous-groupe $H \subset G$ [le corps des invariants \ill{}],
[formé des $g \in G$ tels que] définition de $\Sigma_{y}$ : c'est
$(g \times 1)\Sigma_{y} = \Sigma'_{y}$, ce qui équivaut (comme $K_{0}$
dépend de l'origine \ill{}) à $(g \times 1)\Sigma = \Sigma'$ \ill{} [$H$
doit être indépendant du choix de $y$] : $(g \times 1)(g \times
g^{-1})\Sigma_{y} = \Sigma_{y}$ [vérification : $\Sigma_{y'} = (h \times
h)\Sigma_{y}$ pour $h$ tel que $hy = y'$, donc $(g \times 1)\Sigma_{y'} = (g
\times 1)(h \times h)\Sigma_{y} = (1 \times h)(g \times 1)\Sigma_{y} = (1
\times h)\Sigma'_{y} = \Sigma'_{y'}$], $(g \times 1)(g^{-1} \times
g^{-1})\Sigma = \Sigma'$, et comme $\Sigma = (g^{-1} \times g^{-1})\Sigma$,
cela équivaut [à] $(1 \times g)\Sigma = \Sigma$]. Donc i.e.\ $(1 \times
g^{-1})\Sigma = \Sigma$ \ill{} \ill{}}
\note{tout ce passage est encadré et barré de six longs traits obliques ; il
est très surchargé et n'est rendu qu'en substance}

\emph{NB} Soit $\mathfrak{A}$ une structure CM sur le couple $E, K$, et
considérons $E \hookrightarrow E'$, $K \hookrightarrow K'$ des extensions de
corps. Alors l'idéal $\mathfrak{A}' = \mathfrak{A}(E' \otimes K')$ de $E'
\otimes K'$ est une structure CM de Sh.-Tan.\ sur $E', K'$,
\struck{et réciproquement la théorie de Galois} [$E \otimes K$]
\struck{un idéal CM sur $E, K$} considérons \struck{nous dit} Soit
$\mathfrak{A}$ \struck{structure} un « corps de définition » comme idéal dans
$E_{K}$ ; $\mathfrak{A}$ \ill{} \ill{} invariant par extension $K_{0} \subset
K$. Je dis que celui-ci est invariant par extension $E \to E'$
(vérification triviale). De ceci

\page{148}
résulte qu'il y a un plus petit couple $(E_{0}, K_{0})$ ($E_{0} \subset E$,
$K_{0} \subset K$) tel que $\mathfrak{A}$ provienne d'un idéal de $E_{0}
\otimes K_{0}$. On a vu que $E_{0}, K_{0}$ sont weiliens \ill{}. Si $E_{0}$
est fini sur $\mathbb{Q}$, de degré $2d$, alors $K_{0}$ est fini sur
$\mathbb{Q}$, de degré $\leq \binom{2d}{d}$ [car le s-groupe des $g \in G$
qui invariant une partie à $d$ éléments dans un ens.\ à $2d$ éléments est
d'\struck{\ill{}} indice $\leq \binom{2d}{d}$]. \struck{Si $K$ \ill{} $K$}
Si $E$ fini sur $\mathbb{Q}$ de degré $2d$, alors $\mathfrak{A}(\Sigma) = (E
\otimes K)/\mathfrak{A}$ est fini sur $K$ de degré $d$, et réciproquement.

Soient $E, K$ deux extensions de $\mathbb{Q}$, $t$ un bimodule sur $E, K$,
i.e.\ un module sur $E \otimes K$. On dit que $t$ est du type CM si $t$ est
monogène, et son annulateur est un $\mathfrak{A}$ de type CM. Les bimodules
de type CM sur $E, K$ sont donc (à isom.\ près) \uncertain{paramétrés} par
les structures CM $\Sigma$ sur le couple $E, K$. (si $E \otimes K$ fini sur
$\mathbb{Q}$, i.e.\ $t$ fini sur $E$),
\marginal{$\Sigma$ en surcharge d'un $\mathfrak{A}$}

Un tel $t$ définit [un hom]
\[
  \Psi_{\Sigma} : T_{K} \longrightarrow T_{E}
\]
dont l'action, sur les pts rationnels sur $\mathbb{Q}$ (qui \ill{} \ill{}
idèles) est donnée par
\[
  \Psi_{\Sigma}(x) = \mathrm{d\acute{e}t}_{t/E}(x_{t}) ,
\]
\ill{} fonctions caractéristiques.

Ou si on préfère, en termes de la fonction caractéristique \struck{\ill{}}
$1_{\Sigma}$ de $\Sigma$ (qui est ouvert et fermé), notée
\[
  \Psi_{\Sigma}(x) = N_{E \otimes K/E}(x^{1_{\Sigma}})
\]
et \ill{} de $E$ ; si $K'$ est une extension finie de $K$, [donnant lieu] à
$\Sigma'$ dans $E' \otimes K'$, on trouve commutativité dans

\begin{tikzcd}
T_{K'} \arrow[r, "N_{K'/K}"] \arrow[rrr, bend right=25, "\Psi_{\Sigma'}"'] & T_{K} \arrow[r, "\Psi_{\Sigma}"] & T_{E} \arrow[r, "i_{E'/E}"] & T_{E'}
\end{tikzcd}

\page{149}
Ceci montre que l'on est réduit : étudier [$\Psi_{\Sigma}$ dans le] cas où
$K = K_{0}$, $E = E_{0}$ (donc $E, K$ weiliens et finis).

Identifiant \struck{\ill{}} les duaux de $T_{E}$, $T_{K}$ à valeurs dans
$\overline{\mathbb{Q}}$ [comp.\ \ill{} : $G$)] à $\mathbb{Z}^{I_{E}}$,
$\mathbb{Z}^{I_{K}}$, on constate que l'hom.\ \ill{} \struck{\ill{}} \ill{} de
la \uncertain{conjugaison} défini par $T_{\Sigma}$
\[
  \mathbb{Z}^{I_{E}} \longrightarrow \mathbb{Z}^{I_{K}}
\]
est donné par la matrice $1_{\Sigma}$ considérée comme une fonction
\struck{caractéristique} sur $I_{E} \times I_{K}$). \emph{Donc la
connaissance de $\Sigma$ équivaut à celle de $\Psi_{\Sigma}$} (cette dernière
n'étant d'ailleurs pas quelconque\,…).

Choisissons une clôture algébrique $\overline{E}$ de $E$, \struck{\ill{} $K$}
[$K$ fini (degré $2d'$)] \struck{sur $\mathbb{Q}$}, alors la structure
$\Sigma$ équivaut à la donnée d'un ensemble de $d'$ hom.
\[
  \sigma_{1}, \ldots, \sigma_{d'} : K \longrightarrow \overline{E}
\]
tels que pour tte conj.\ $\tau$ de $G = \mathrm{Gal}(\overline{E} |
\mathbb{Q})$, \struck{\ill{}} les $\{\tau\sigma_{i}, \sigma_{i}\}$ forment
l'ens.\ de tous les hom $K \to \overline{E}$.

On aura alors
\[
  \Psi_{\Sigma}(x) = \sigma_{1}(x)\sigma_{2}(x) \cdots \sigma_{d'}(x)
\]
\struck{les $\sigma_{i} | K_{r}$ forment \ill{} $x \in K_{r}$ \ill{}}.
Si $K$ est weilien, \struck{\ill{}} [et on a] les hom $K_{r} \to
\overline{E}$, et on \ill{} \struck{les hom \ill{} $\tau\sigma_{i}(x) =
\sigma_{i}(x)$, donc \ill{} $\Psi_{\Sigma}(x) = N_{K}$}
\[
  \Psi_{\Sigma} | T_{K_{r}} = N_{K_{r}/\mathbb{Q}}
\]
\note{au-dessus, une première formule biffée ; le $T$ de $T_{K_{r}}$ est
écrit en surcharge}
Par suite, $\Psi_{\Sigma}$ s'annule sur $V_{K_{r}} \subset T_{K_{r}} \subset
T_{K}$. De ceci, et de \ill{}, on conclut que \emph{si $K$ fini sur
$\mathbb{Q}$}, alors $\Psi_{\Sigma} : T_{K} \to T_{E}$ se \emph{factorise
par} $S^{\circ}_{K} \to T_{E}$ :

\page{150}
\[
  T_{K} \xrightarrow{\;\mathrm{can}\;} S^{\circ}_{K}
  \xrightarrow{\;\Psi_{\Sigma}\;} T_{E}
\]
\marginal{au-dessus de $\Psi_{\Sigma}$ : « pas changé de notation »}

Regardons-nous de l'homomorphisme
\[
  i_{K} : \mathbb{G}_{m,\mathbb{Q}} \longrightarrow S^{\circ}_{K}
\]
Je dis qu'on a
\[
  \boxed{\Psi_{\Sigma}\, i_{K} = i_{E/\mathbb{Q}}}
\]
\marginal{a)}
Comme \ill{}
\[
  i_{K}(x)^{d'} = i_{K/\mathbb{Q}}(x) \qquad (x \in \mathbb{Q}^{*})
\]
\note{un signe biffé entre $i_{K}$ et $(x)$}
la relation est équivalente (\ill{} \ill{} : \ill{} d'\ill{}) à
\[
  \Psi_{\Sigma}(x) = x^{d'} \qquad \text{pour } x \in \mathbb{Q}^{*}
\]
qui est trivialement vérifiée.

[Quelle est l'image de $\Psi_{\Sigma}$ ? Je \ill{} \ill{} de
$S^{\circ}_{K}$, désignons l'image de $V'_{K}$, le ss-\struck{\ill{}}
\uncertain{tore} compact maximal de $S^{\circ}_{K}$, on
\[
  \Psi_{\Sigma}(S^{\circ}_{K}) = \Psi_{\Sigma}(V'_{K}) \cdot
  i_{E/\mathbb{Q}}(\mathbb{G}_{m})
\]
On va bien voir \struck{\ill{}} [\uncertain{brancher}] au cas $E =
E_{0}$\,….]

\emph{Question} Pour $K$ fixé, [tel que $S^{\circ}_{K} \to \mathbb{G}_{m}$
\ill{} \ill{} isom.,] les hom $\Psi_{\Sigma} : S^{\circ}_{K} \to T_{E}$
associés aux diverses structures CM sur $K$ \struck{sont} [\ill{} $K$
weilien \ill{}] forment-ils une famille \emph{fidèle} ?

[Soit $\Sigma = \{\sigma_{1}, \ldots, \sigma_{d'}\}$ une structure CM sur
$(K, \overline{\mathbb{Q}})$. \struck{Alors si $x \in K^{*}$ \ill{} pour tt
$1 \leq i \leq d'$,} $\Sigma_{i}$ la structure CM déduite de $\Sigma$ en
remplaçant $\sigma_{i}$ par $\bar{\sigma}_{i} = \tau\sigma_{i}$. Alors, si
$x \in K^{*}$ est tel que $\Psi_{\Sigma}(x) = 1$ et $\Psi_{\Sigma_{i}}(x) =
1$, on aura en quotient $\sigma_{i}(x) = \bar{\sigma}_{i}(x)$, donc on aura
$x \in K_{r}$. De plus, si ceci est vrai pour tt $x$, on aura
$N_{K_{r}/\mathbb{Q}}(x) = 1$, donc $x \in V'_{K_{r}}$.
\struck{Donc} $\Psi_{\Sigma}(x) = 1$ \uncertain{revient} donc à : l'image de
$x$ dans $S^{\circ}_{K}$ est $1$. Cet argument \uncertain{marche} aussi en
remplaçant $\mathbb{Q}$ par une alg.\ quelconque sur $\mathbb{Q}$
\struck{[car $\bigcap \mathrm{Ker}(\sigma_{i}, \bar{\sigma}_{i}) = K_{r}$ est
stable par]}
\note{la ligne biffée est reprise en tête du feuillet suivant}

\page{151}
(\emph{NB} $K_{r} = \bigcap_{i} \mathrm{Ker}(\sigma_{i}, \bar{\sigma}_{i} :
K \rightrightarrows \overline{\mathbb{Q}})$ est stable par tt changement de
base $\mathbb{Q} \to \Lambda$). Donc
\[
  \mathrm{Ker}\,\Psi_{\Sigma} \cap \bigcap_{1 \leq i \leq d'}
  \mathrm{Ker}\,\Psi_{\Sigma_{i}} = 1 \qquad \text{dans } S^{\circ}_{K}
\]
\note{« Ker » de deux flèches désigne ici leur noyau au sens du
conoyau-égalisateur : le sous-corps où $\sigma_{i}$ et $\bar{\sigma}_{i}$
coïncident}

\emph{Question} Regardons les hom.\ $\Psi_{\Sigma} : S^{\circ}_{K} \to T_{E}$,
avec $E$ \emph{quadratique}. Y en a-t-il assez, du moins en \uncertain{filtrant}
\ill{} sur $K$ croissants ? Or un $E$ a, relativement à une \ill{}
extension $K$ de $\mathbb{Q}$, \struck{\ill{}} $0$ structure CM si $E$ n'est
pas \ill{} [(\ill{} \ill{}) \ill{} \ill{} \ill{}] dans $K$, et deux
\uncertain{exactement} \ill{} \ill{} dans $K$. L'hom $\Psi_{K}$
correspondant \uncertain{est} \struck{Tr} $N_{K/E}$.

Donc pour $K$ fixé, les hom $\Psi_{\Sigma} : S^{\circ}_{K} \to T_{E}$
correspondant, par paires conjuguées, aux plongements $E \to K$ ; donc pour
voir s'il y a « assez » de \ill{} il faut voir s'il y a assez [(supposé
weilien), \ill{}] de sous-corps quadratiques de $K$. Si $K$ [\ill{} \ill{}
\ill{}] \ill{} \ill{} sous-corps quadratiques i.e.\ si \struck{\ill{}} [$K$
\ill{}] \ill{} de groupe de Galois $(\mathbb{Z}/2\mathbb{Z})^{N}$, dans le
noyau de
\[
  N_{K/L} : S^{\circ}_{K} \longrightarrow S^{\circ}_{L}
\]
est \struck{\ill{}} [un] tore de \ill{}, qui est dans le noyau de tous les
$\Psi_{K}$ envisagés.
\note{les deux tiers inférieurs de la page sont d'une main très rapide ; la
phrase finale est peu sûre}

\page{153}
\note{nouveau feuillet, qui ne continue pas le précédent}
$K$ corps CM [galoisien] sur $\mathbb{Q}$, groupe de Galois $G$, $\gamma =
\{1, \rho\} \simeq \mathbb{Z}/2\mathbb{Z} \subset G$, $\rho$ la conj.\
complexe.

$\overline{\mathbb{Q}}$ [une] clôture algébrique de $\mathbb{Q}$
\struck{\ill{} $K$}.

\struck{$I$ ensemble des plongements de $K$ dans $\overline{\mathbb{Q}}$,
$G$ y opère par transport de structure $g.\sigma = \sigma \circ g^{-1}$, de
façon simplement transitive.}

$M$ groupe des fonctions \struck{\ill{}} $\varphi : G \to \mathbb{Z}$ telles
que
\[
  \varphi(\rho\sigma) = -\varphi(\sigma) \qquad \text{pour tt } \sigma \in G
\]
C'est un $G$-module, qui définit [\uncertain{par descente et dualité}] un
tore $T$ sur $\mathbb{Q}$.
\marginal{\emph{NB} $M \simeq \check{M}$ comme $G$-module\struck{s}.}

(1) On veut calculer le sous-groupe [$H^{1}(\mathbb{Q}, T)'$] de
$H^{1}(\mathbb{Q}, T)$ formé des classes qui sont nulles en tout
\struck{place} \uncertain{place finie}. On trouve par \struck{\ill{}
$P = I/\gamma$}, \ill{}, une \struck{\ill{}} [suite exacte] canonique
\[
  0 \longrightarrow H^{1}(\mathbb{Q}, T)' \xrightarrow{\;\alpha\;}
  \hat{H}_{0}(G, Y_{\infty} \otimes_{\mathbb{Z}} \check{M})
  \xrightarrow{\;\beta\;} \hat{H}_{0}(G, \check{M}) \longrightarrow 0
\]
(où $Y_{\infty}$ est le gpe libre engendré par les places \uncertain{à
l'infini} de $K$, \ill{} \ill{}, $\beta$ est déduit de l'augmentation
naturelle $Y_{\infty} \to \mathbb{Z}$ (somme des valeurs)).

Si on \uncertain{prend} $\overline{\mathbb{Q}}$ la clôture algébrique de
$\mathbb{Q}$ dans $\mathbb{C}$, \ill{} $\mathbb{Z}'$ est le $\gamma$-module
$\mathbb{Z}$ où $\rho$ opère par $-\mathrm{id}_{\mathbb{Z}}$. On a
$M \simeq \mathbb{Z}' \otimes_{\gamma} G$, \ill{}
\note{dans $\mathbb{Z}' \otimes_{\gamma} G$, un premier $\mathbb{Z}$ est
biffé}

\begin{tikzcd}[column sep=small]
0 \arrow[r] & H^{1}(\mathbb{Q}, T)' \arrow[r] & \hat{H}_{0}(\gamma, Y'_{\infty}) \arrow[r] \arrow[d, "\simeq"] & \hat{H}_{0}(\gamma, \mathbb{Z}') \arrow[r] \arrow[d, no head] & 0 \\
 & & Y_{\infty} \otimes \mathbb{Z}/2\mathbb{Z} \arrow[r, two heads] & \mathbb{Z}/2\mathbb{Z} &
\end{tikzcd}

\note{les deux traits verticaux sont, sur la page, un $\simeq$ couché et un
signe d'égalité ; après le dernier $\mathbb{Z}/2\mathbb{Z}$, un mot biffé}

et donc on trouve
\[
  H^{1}(\mathbb{Q}, T)' \simeq \text{gpe des combinaisons linéaires}
\]
[des] places à l'infini de $K$, à somme des coefficients nulle.
\marginal{i.e.\ coeff.\ dans $\mathbb{F}_{2}$}

\page{154}
\note{feuillet isolé, d'un autre sujet (familles épimorphiques dans une
catégorie), barré de quatre longs traits obliques au crayon}
Lorsque \uncertain{cependant} dans $C$ les produits fibrés existent, alors il
n'y a pas de différence entre épim.\ eff.\ et épim.\ strict, il n'y a plus
que quatre (\uncertain{au lieu de} \uncertain{cinq}) variantes
d'épimorphicité pour des familles, et ces \struck{variantes} propriétés se
reconnaissent déjà sur $R$ le crible $R$ (cet \ill{} \ill{} \ill{}).

\emph{Prop.\ 0}. Les \struck{cinq} notions \struck{envisagées}
[« épim. », « épim.\ stricte » \struck{\ill{}}] sont \struck{la propriété}
stables par \ill{} [les trois dernières variantes], \ill{} de la famille,
\ill{} ne dépendent que des \ill{} quarrables.

Soit $G$ une partie de $\mathrm{Ob}\, C$, $\underline{G}$ la sous-catégorie
pleine engendrée. \struck{\ill{}}

\emph{Prop.\ 1}. Conditions équivalentes

a) Pour tout $X \in \mathrm{Ob}\, C$ la famille des flèches de but $X$,
source $\in G$, est épimorphique (resp.\ épimorphique stricte).

b) Pour tout $X \in \mathrm{Ob}\, C$, considérant $\underline{G}_{/X}
\xrightarrow{\;i\;} C$ et le système inductif évident
\[
  \varinjlim_{T \in \underline{G}/X} T \longrightarrow X ,
\]
[i.e.\ un \ill{} en sens inverse des foncteurs covariants [cov.]\ en $Y$],
cet système est épimorphique \struck{\ill{}} (resp.\ isomorphique) [i.e.\
l'hom.\ \ill{} de foncteurs cov.\ en $Y$ est un mono, resp.\ un isom].

\emph{Dém.} Le cas non respé est une tautologie, regardons le cas respé. La
condition a) signifie que, si $u_{i} : G_{i} \to X$ sont les différents
objets de $\underline{G}/X$, alors $\forall Y$ (i) $h_{Y}(X) \to \prod_{i}
h_{Y}(G_{i})$ injectif et (ii) Pour des $\xi_{i} \in h_{Y}(G_{i})$, tels
\note{la page s'arrête sur ce mot ; la suite n'est pas dans ce lot (le
feuillet 155 est blanc)}

\section{Corps de classes commutatif = Motifs à groupe fondamental motivique [connexe] commutatif (= Motifs métabéliens)}

\note{mots de sa main, seuls sur une couverture brune (feuillet 156) ; les
crochets sont les siens, et « connexe » est une lecture douteuse. Les
feuillets 158 et 160, d'une encre noire sur un papier jauni, suivent}

\page{158}
$K$ corps de nbres.

$A$ anneau des entiers, \ill{} fini libre sur $\mathbb{Z}$.

$\underline{\underline{A}}$ schéma en anneaux sur $\mathbb{Z}$

$\underline{\underline{A}}^{*}$ schéma en groupes.

$\underline{\underline{A}}^{*}(\mathbb{Z}) = \emph{unités} = E$.

$E^{\circ} \subset E$ \emph{unités} \emph{\uncertain{positives}}, d'indice
fini de la forme $2^{\alpha}$ ($\alpha \leq r_{1}$) dans $E$, égal à $E$ si
$K$ est purement imaginaire.

\struck{\ill{}} Le groupe de Galois de $K$ rendu abélien est [une extension
de $\mathrm{Pic}(A)$ par] alors
\[
  \underline{\underline{A}}^{*}(\hat{\mathbb{Z}})/E^{\circ}
\]
ou encore par
\[
  \underline{\underline{A}}^{*}(\struck{\ill{}}\, \mathbb{Z} \times
  \mathbb{R}^{*}) / \struck{\text{composante}}\,
  \underline{\underline{A}}^{*}(\mathbb{Z})
\]
divisé par la composante connexe.
\marginal{vérification en \ill{}}

Le schéma en groupes [connexe] sur $\mathbb{Z}$ associé à $K$ par la
classification des motifs abéliens sur $K$, semble être alors
\[
  G^{\circ} = \underline{\underline{A}}^{*} / \mathfrak{N} ,
\]
où $\mathfrak{N}$ est \struck{\ill{}} le sous-schéma en groupe de
$\underline{\underline{A}}^{*}$ \emph{adhérence} de \ill{} l'un des
\struck{\ill{}} [sections] $E^{\circ}$ sur $\mathbb{Z}$. + \ill{}
\struck{schéma} de Zariski fermé et \ill{} \ill{} de $E^{\circ}$ dans
$\underline{K}^{*}$.
\note{les quatre dernières lignes sont peu sûres ; le soulignement double
de $\underline{K}^{*}$ est le sien}

\page{160}
\struck{\ill{}} [Extension canonique de $\gamma$ par (le tore) $G^{\circ}$],

\begin{tikzcd}
0 \arrow[r] & G^{\circ} \arrow[r] & G \arrow[r] & \gamma \arrow[r] \arrow[d, no head] & 0 \\
 & & & \mathrm{Pic}(A)^{\mathrm{abs}} &
\end{tikzcd}

\note{le trait vertical est un signe d'égalité ; l'exposant de
$\mathrm{Pic}(A)$ se lit « abs », lecture douteuse}

Ceci posé, à toute place $\xi$ de $K$, sont associés des foncteurs

\begin{tikzcd}[column sep=small, row sep=small, nodes={font=\scriptsize}]
M_{\circ} \arrow[r, "\tilde{\xi}"] \arrow[rr, bend right=20, "\text{fibre en } \xi"'] & \text{Vectoriels sur } \mathbb{Q} \text{ avec opération de } G & \text{réseaux de Hodge}
\end{tikzcd}

\note{sous $M_{\circ}$, un signe d'égalité et « motifs abéliens » ; sous la
flèche $\tilde{\xi}$, un $\sim$ ; au-dessus du but, « fibre en $\xi$ » ;
« réseaux » est une lecture douteuse. Après « réseaux de Hodge », entre
parenthèses : « + structure réelle si la place est réelle »}

d'où un foncteur canonique
\[
  \text{Vectoriels sur } \mathbb{Q} \text{ avec opération de } G
  \xrightarrow{\;\Phi_{\xi}\;} \text{réseaux de Hodge}
\]
(+ structure réelle si $\xi$ réelle).

\struck{En $\Phi_{\xi}$} Ainsi, si $V$ est un \struck{\ill{}}
$\mathbb{Q}$-vectoriel sur lequel $G$ opère, alors pour toute place $\xi$ de
$K$, on trouve \emph{une structure de Hodge sur $V \otimes_{\mathbb{Q}}
\mathbb{C}$}. Il suffit d'ailleurs d'avoir les \uncertain{opérations} de
$G^{\circ}$ sur $V$, i.e.\ de $\underline{K}^{*}$ sur $V$ [trivial sur $E$],
+ une place de [structure de] $\tilde{\xi}$ de $\widetilde{K}$, pour en
conclure la \uncertain{structure} de Hodge.
\marginal{$\widetilde{K}$ : corps de classes absolu}

\end{document}
